% ============================================================
%  CHAPTER 7: IMPLEMENTATION DETAILS
% ============================================================

\chapter{IMPLEMENTATION DETAILS}
\label{chap:implementation}

\section{Project Directory Structure and Repository Organization}
\label{sec:directory_structure}

The QuantumSafe VPN codebase is organized into two primary sub-projects within a unified engineering repository: the native Kotlin Android client application and the companion multi-client Python/Linux gateway server.

\begin{lstlisting}[basicstyle=\ttfamily\scriptsize, caption={Repository Directory Structure.}]
VPN_APP/
|-- android/                             # Native Android Client Project
|   |-- app/
|   |   |-- src/main/
|   |   |   |-- AndroidManifest.xml      # OS permissions & service declarations
|   |   |   |-- java/com/quantumsafe/vpn/
|   |   |   |   |-- MainActivity.kt      # Application entry point & navigation
|   |   |   |   |-- QuantumSafeVpnService.kt # Core VpnService & TUN packet loop
|   |   |   |   |-- crypto/
|   |   |   |   |   `-- PqcEngine.kt     # Pure-Kotlin PQC & AES-GCM engine
|   |   |   |   |-- model/
|   |   |   |   |   `-- VpnModels.kt     # Data models & state representations
|   |   |   |   |-- service/
|   |   |   |   |   `-- VpnServiceBridge.kt # Inter-process communication bridge
|   |   |   |   |-- storage/
|   |   |   |   |   `-- SecurityManager.kt  # Android Keystore configuration storage
|   |   |   |   |-- ui/
|   |   |   |   |   |-- MainScreen.kt    # Root UI scaffold & navigation bar
|   |   |   |   |   |-- screens/
|   |   |   |   |   |   |-- HomeScreen.kt # Primary connect & telemetry dashboard
|   |   |   |   |   |   |-- ServerConfigScreen.kt # Server endpoint configuration
|   |   |   |   |   |   |-- SecurityDetailsScreen.kt # Cryptographic telemetry
|   |   |   |   |   |   `-- SettingsScreen.kt     # Diagnostics & terminal logs
|   |   |   |   |   `-- theme/
|   |   |   |   |       `-- Theme.kt     # Minimal white design tokens
|   |   |   |   `-- viewmodel/
|   |   |   |       `-- VpnViewModel.kt  # Reactive MVVM business logic
|   |   |-- build.gradle                 # Module dependencies & compiler flags
|   `-- gradle/                          # Gradle wrapper binaries
|-- vpn/                                 # Linux Gateway Server Daemon
|   |-- tunnel_server.py                 # Multi-client TCP listener & dispatch
|   |-- ip_pool.py                       # Thread-safe 10.8.0.0/24 IP pool
|   |-- nat_manager.py                   # Automated iptables NAT masquerading
|   |-- tun.py                           # Linux /dev/net/tun interface wrapper
|   `-- config.py                        # Server configuration loader
|-- pqc/                                 # Server Post-Quantum Modules
|   |-- mlkem.py                         # ML-KEM-1024 encapsulation & decapsulation
|   `-- mldsa.py                         # ML-DSA-65 identity generation & signing
|-- crypto/                              # Classical Symmetric Primitives
|   `-- classical.py                     # HKDF-SHA256 & AES-256-GCM AEAD
|-- tests/                               # Comprehensive Automated Test Suites
|-- dist/                                # Built Production Debug & Release APKs
`-- releases/                            # Distribution Artifacts & Checksums
\end{lstlisting}

\section{Android Gradle Build Configuration}
\label{sec:gradle_build_config}

The Android mobile client is configured via Gradle 8.9 utilizing Android Gradle Plugin (AGP) version 8.3.2 and Kotlin 1.9.22.

The module build script (\texttt{android/app/build.gradle}) defines strict SDK targets:
\begin{lstlisting}[language=Kotlin, caption={Gradle Build Configuration Excerpt.}]
android {
    namespace = "com.quantumsafe.vpn"
    compileSdk = 34

    defaultConfig {
        applicationId = "com.quantumsafe.vpn"
        minSdk = 26              // Android 8.0 (Oreo)
        targetSdk = 34           // Android 14 (Upside Down Cake)
        versionCode = 1
        versionName = "1.0.0"
        testInstrumentationRunner = "androidx.test.runner.AndroidJUnitRunner"
    }

    buildFeatures {
        compose = true
    }

    composeOptions {
        kotlinCompilerExtensionVersion = "1.5.8"
    }

    compileOptions {
        sourceCompatibility = JavaVersion.VERSION_17
        targetCompatibility = JavaVersion.VERSION_17
    }

    kotlinOptions {
        jvmTarget = "17"
        freeCompilerArgs += listOf(
            "-opt-in=kotlinx.coroutines.ExperimentalCoroutinesApi",
            "-opt-in=androidx.compose.material3.ExperimentalMaterial3Api"
        )
    }
}
\end{lstlisting}

Key dependencies include:
\begin{itemize}
    \item \texttt{androidx.compose.material3:material3:1.2.0}: Material 3 composables.
    \item \texttt{androidx.lifecycle:lifecycle-viewmodel-compose:2.7.0}: Reactive ViewModel binding.
    \item \texttt{org.jetbrains.kotlinx:kotlinx-coroutines-android:1.7.3}: Asynchronous coroutine workers.
    \item \texttt{androidx.security:security-crypto:1.1.0-alpha06}: Hardware-backed Keystore encryption.
\end{itemize}

\section{Android Manifest and Foreground Service Compliance}
\label{sec:manifest_service_declaration}

Android~14 (API Level~34) enforces strict security constraints on background execution, requiring explicit foreground service type declarations to prevent silent application suspension.

In \texttt{AndroidManifest.xml}:
\begin{lstlisting}[language=XML, caption={AndroidManifest Service and Permission Declarations.}]
<manifest xmlns:android="http://schemas.android.com/apk/res/android">

    <uses-permission android:name="android.permission.INTERNET" />
    <uses-permission android:name="android.permission.ACCESS_NETWORK_STATE" />
    <uses-permission android:name="android.permission.FOREGROUND_SERVICE" />
    <uses-permission android:name="android.permission.FOREGROUND_SERVICE_SYSTEM_EXEMPTED" />
    <uses-permission android:name="android.permission.POST_NOTIFICATIONS" />

    <application
        android:label="QuantumSafe VPN"
        android:icon="@mipmap/ic_launcher"
        android:theme="@style/Theme.QuantumSafeVPN">

        <activity
            android:name=".MainActivity"
            android:exported="true"
            android:configChanges="orientation|screenSize">
            <intent-filter>
                <action android:name="android.intent.action.MAIN" />
                <category android:name="android.intent.category.LAUNCHER" />
            </intent-filter>
        </activity>

        <service
            android:name=".QuantumSafeVpnService"
            android:permission="android.permission.BIND_VPN_SERVICE"
            android:foregroundServiceType="systemExempted"
            android:exported="false">
            <intent-filter>
                <action android:name="android.net.VpnService" />
            </intent-filter>
        </service>
    </application>
</manifest>
\end{lstlisting}

The service declares \texttt{android:foregroundServiceType="systemExempted"} in accordance with Android 14 guidelines for VPN providers, ensuring that the operating system's battery optimization daemon does not terminate the socket during continuous background operation.

\section{Core Android VPN Service Implementation}
\label{sec:core_vpn_service_impl}

The \texttt{QuantumSafeVpnService} class extends \texttt{android.net.VpnService}, managing the core lifecycle of the tunnel connection.

\subsection{Foreground Notification Lifecycle}
When initiated, the service builds and binds an ongoing system notification displaying connection state, target gateway IP, and an active ``Disconnect'' pending intent:
\begin{lstlisting}[language=Kotlin, caption={Foreground Service Notification Binding.}]
private fun createNotification(statusText: String): Notification {
    val channelId = "quantumsafe_vpn_channel"
    val channel = NotificationChannel(
        channelId,
        "VPN Connection Service",
        NotificationManager.IMPORTANCE_LOW
    )
    getSystemService(NotificationManager::class.java).createNotificationChannel(channel)

    val disconnectIntent = PendingIntent.getService(
        this, 0,
        Intent(this, QuantumSafeVpnService::class.java).apply { action = "ACTION_DISCONNECT" },
        PendingIntent.FLAG_IMMUTABLE
    )

    return NotificationCompat.Builder(this, channelId)
        .setContentTitle("QuantumSafe VPN")
        .setContentText(statusText)
        .setSmallIcon(R.drawable.ic_vpn_lock)
        .setOngoing(true)
        .addAction(R.drawable.ic_stop, "Disconnect", disconnectIntent)
        .build()
}
\end{lstlisting}

\subsection{Socket Bypassing and Non-blocking Connection}
Prior to establishing the TCP connection to the gateway on port 7000, the socket must be protected:
\begin{lstlisting}[language=Kotlin, caption={Socket Protection in QuantumSafeVpnService.}]
val socket = Socket()
val isProtected = protect(socket)
if (!isProtected) {
    Log.e(TAG, "Critical failure: Socket protect returned false!")
    broadcastError("Failed to protect transport socket from VPN loop.")
    return
}
socket.tcpNoDelay = true
socket.connect(InetSocketAddress(serverHost, serverPort), 10000)
\end{lstlisting}
Enabling \texttt{tcpNoDelay = true} disables Nagle's algorithm, ensuring that small encapsulated packets (such as interactive SSH keystrokes or TCP ACKs) are transmitted immediately without artificial 40\,ms buffering delays.

\section{Pure-Kotlin Post-Quantum Cryptographic Engine (\texttt{PqcEngine.kt})}
\label{sec:pqc_engine_impl}

The cryptographic engine \texttt{PqcEngine.kt} provides high-performance, memory-safe execution of lattice arithmetic, hashing, and authenticated encryption.

\subsection{Handshake Execution Logic}
The client executes the 4-stage post-quantum cryptographic exchange over the connected socket:
\begin{lstlisting}[language=Kotlin, caption={Handshake Execution in PqcEngine.}]
fun performHandshake(socket: Socket): HandshakeResult {
    val inputStream = DataInputStream(socket.getInputStream())
    val outputStream = DataOutputStream(socket.getOutputStream())

    // Stage 1: Receive ML-DSA-65 Server Public Key (1952 Bytes)
    val serverDsaPk = readFrame(inputStream)
    if (serverDsaPk.size != 1952) {
        throw SecurityException("Invalid server ML-DSA public key size: ${serverDsaPk.size}")
    }

    // Stage 2: Generate ML-KEM-1024 Ephemeral Keypair & Send Public Key (1568 Bytes)
    val (kemPk, kemSk) = generateMlkemKeypair()
    sendFrame(outputStream, kemPk)

    // Stage 3: Receive Ciphertext (1568B) and ML-DSA Signature (3309B)
    val ciphertext = readFrame(inputStream)
    val signature = readFrame(inputStream)

    // Verify Server Signature over Transcript
    val transcript = "QuantumSafeVPN-v1".toByteArray(Charsets.UTF_8) + kemPk + ciphertext
    val isSigValid = verifyMldsa(serverDsaPk, transcript, signature)
    if (!isSigValid) {
        throw SecurityException("Server identity verification failed! Signature invalid.")
    }

    // Decapsulate Shared Secret
    val sharedSecret = decapsulateMlkem(kemSk, ciphertext)

    // Stage 4: Receive Configuration JSON Payload
    val configJsonBytes = readFrame(inputStream)
    val config = parseConfigJson(configJsonBytes)

    // Derive AES-256-GCM Session Key via HKDF-SHA256
    val sessionKey = deriveSessionKey(sharedSecret)

    return HandshakeResult(sessionKey, config)
}
\end{lstlisting}

\subsection{Symmetric Packet Framing and Galois/Counter Mode Encryption}
For line-rate packet transmission, \texttt{PqcEngine} encapsulates data frames:
\begin{lstlisting}[language=Kotlin, caption={AES-256-GCM Packet Framing in PqcEngine.}]
fun encryptPacket(plaintext: ByteArray, sequenceNumber: Long): ByteArray {
    val iv = ByteArray(12)
    secureRandom.nextBytes(iv)

    // Prepend 64-bit monotonic sequence number
    val payloadBuffer = ByteBuffer.allocate(8 + plaintext.size)
    payloadBuffer.order(ByteOrder.BIG_ENDIAN)
    payloadBuffer.putLong(sequenceNumber)
    payloadBuffer.put(plaintext)
    val dataToEncrypt = payloadBuffer.array()

    // Initialize AES-256-GCM Cipher
    val cipher = Cipher.getInstance("AES/GCM/NoPadding")
    val spec = GCMParameterSpec(128, iv) // 128-bit authentication tag
    cipher.init(Cipher.ENCRYPT_MODE, secretKey, spec)
    val ciphertextAndTag = cipher.doFinal(dataToEncrypt)

    // Assemble final wire frame: [IV (12B)] + [Ciphertext + Tag (16B)]
    val frameBuffer = ByteBuffer.allocate(12 + ciphertextAndTag.size)
    frameBuffer.put(iv)
    frameBuffer.put(ciphertextAndTag)
    return frameBuffer.array()
}
\end{lstlisting}

\section{Pure-Kotlin Keccak-$f[1600]$ SHA3-256 Implementation}
\label{sec:keccak_impl_walkthrough}

To eliminate cryptographic divergence with Python's \texttt{hashlib.sha3\_256}, \texttt{PqcEngine} implements a pure-Kotlin Keccak sponge function conforming to NIST FIPS 202:

\begin{lstlisting}[language=Kotlin, caption={Pure-Kotlin Keccak-f[1600] Permutation Round.}]
private fun keccakF1600(state: LongArray) {
    for (round in 0 until 24) {
        // Theta step
        val c = LongArray(5)
        for (x in 0 until 5) {
            c[x] = state[x] xor state[x + 5] xor state[x + 10] xor state[x + 15] xor state[x + 20]
        }
        for (x in 0 until 5) {
            val d = c[(x + 4) % 5] xor java.lang.Long.rotateLeft(c[(x + 1) % 5], 1)
            for (y in 0 until 5) {
                state[x + 5 * y] = state[x + 5 * y] xor d
            }
        }

        // Rho and Pi steps
        val b = LongArray(25)
        for (x in 0 until 5) {
            for (y in 0 until 5) {
                b[y + 5 * ((2 * x + 3 * y) % 5)] = 
                    java.lang.Long.rotateLeft(state[x + 5 * y], KECCAK_ROTATION_CONSTANTS[x][y])
            }
        }

        // Chi step
        for (x in 0 until 5) {
            for (y in 0 until 5) {
                state[x + 5 * y] = b[x + 5 * y] xor (b[(x + 1) % 5 + 5 * y].inv() and b[(x + 2) % 5 + 5 * y])
            }
        }

        // Iota step
        state[0] = state[0] xor KECCAK_ROUND_CONSTANTS[round]
    }
}
\end{lstlisting}

\section{Asynchronous Bidirectional Packet Forwarding Loop}
\label{sec:packet_forwarding_loop}

Once the virtual interface is established, \texttt{QuantumSafeVpnService} spawns two asynchronous coroutine workers running on \texttt{Dispatchers.IO}:

\begin{lstlisting}[language=Kotlin, caption={Uplink Packet Processing Worker Coroutine.}]
private fun CoroutineScope.launchUplinkWorker(
    pfd: ParcelFileDescriptor,
    socket: Socket,
    pqcEngine: PqcEngine
) = launch(Dispatchers.IO) {
    val inputStream = FileInputStream(pfd.fileDescriptor)
    val outputStream = DataOutputStream(socket.getOutputStream())
    val packetBuffer = ByteArray(32767)

    try {
        while (isActive && isRunning) {
            val bytesRead = inputStream.read(packetBuffer)
            if (bytesRead > 0) {
                val rawIpPacket = packetBuffer.copyOf(bytesRead)
                val seq = sendSequenceCounter.incrementAndGet()

                val encryptedFrame = pqcEngine.encryptPacket(rawIpPacket, seq)

                // Write framed packet: [4B Length] + [Encrypted Data]
                synchronized(outputStream) {
                    outputStream.writeInt(encryptedFrame.size)
                    outputStream.write(encryptedFrame)
                    outputStream.flush()
                }

                // Update Telemetry
                totalBytesSent.addAndGet(encryptedFrame.size.toLong())
            }
        }
    } catch (e: Exception) {
        if (isRunning) handleTunnelError("Uplink read failure: ${e.message}")
    }
}
\end{lstlisting}

\section{Python Linux Gateway Daemon Implementation}
\label{sec:python_gateway_impl}

On the server side, \texttt{vpn/tunnel\_server.py} manages incoming client connections using Python's \texttt{threading} and \texttt{socket} modules:

\subsection{Multi-Client IP Allocation and Dispatch}
When a client connects, the server performs the handshake, allocates a virtual IP from \texttt{ip\_pool.py}, and spawns a dedicated client handler:
\begin{lstlisting}[language=Python, caption={Python Server Multi-Client Handshake and Registration.}]
def handle_client_connection(tun, conn, addr):
    try:
        session_key, assigned_ip = perform_handshake(conn, addr)
        session_router.register_client(assigned_ip, conn, session_key)
        print(f"[+] Client {addr} active at virtual IP {assigned_ip}")

        last_seq = 0
        while True:
            encrypted_frame = recv_frame(conn)
            plaintext = decrypt(session_key, encrypted_frame)

            # Extract 64-bit sequence counter
            seq = struct.unpack("!Q", plaintext[:8])[0]
            if seq <= last_seq:
                print(f"[!] Warning: Dropped replay packet (seq={seq}, last={last_seq})")
                continue
            last_seq = seq

            raw_ip = plaintext[8:]
            tun.write_packet(raw_ip) # Inject to Linux tun0

    except Exception as e:
        print(f"[-] Client {addr} disconnected: {e}")
    finally:
        session_router.unregister_client(assigned_ip)
        ip_pool.release(assigned_ip)
        conn.close()
\end{lstlisting}

\subsection{TUN Dispatcher Thread}
A dedicated singleton background thread continuously reads return packets from \texttt{tun0}, parses the destination IP from the IPv4 header, and dispatches the packet to the appropriate client socket:
\begin{lstlisting}[language=Python, caption={TUN Dispatcher Loop in tunnel\_server.py.}]
def tun_dispatcher_loop(tun):
    while True:
        packet = tun.read_packet()
        if not packet:
            continue
        dest_ip = SessionRouter.extract_dest_ip(packet)
        client_info = session_router.get_client(dest_ip)

        if client_info:
            sock = client_info["socket"]
            key = client_info["session_key"]
            client_info["send_seq"] += 1
            seq_packet = struct.pack("!Q", client_info["send_seq"]) + packet
            encrypted = encrypt(key, seq_packet)
            send_frame(sock, encrypted)
\end{lstlisting}

\section{Build and Packaging Process}
\label{sec:build_packaging_process}

The Android mobile client is compiled using Gradle:
\begin{lstlisting}[language=bash, caption={APK Compilation and Packaging Commands.}]
# Clean and assemble debug APK
./gradlew assembleDebug

# Assemble release APK
./gradlew assembleRelease

# Align and sign release APK
zipalign -v -p 4 app-release-unsigned.apk app-release-aligned.apk
apksigner sign --ks release.keystore --ks-key-alias quantumsafe \
    --out QuantumSafe-VPN-release.apk app-release-aligned.apk
\end{lstlisting}

The resulting production artifacts reside in \texttt{dist/}:
\begin{itemize}
    \item \texttt{QuantumSafe-VPN-debug.apk} ($19.0$\,MB): Includes debugging symbols and logging.
    \item \texttt{QuantumSafe-VPN-release.apk} ($14.1$\,MB): Optimized with R8 code shrinking and resource minification.
\end{itemize}
SHA-256 verification hashes are generated and validated in \texttt{dist/SHA256SUMS.txt}.
